\section{Cadena probatoria continua de la publicación excepcional: del despliegue APP a Paley--Witt,
\texorpdfstring{\(V^\natural\)}{V natural} y Moonshine}
\label{sec:prueba-corredor-excepcional}

La publicación excepcional parte del lector de incidencia
\(\mathfrak E_{\rm exc}\) de \eqref{eq:frontal-doble-proyeccion}. No toma como
entrada un código, un retículo o un grupo excepcional ya dados: toma la
frontera, la orientación, la hoja y la memoria de una historia enriquecida y
las expresa en el calibre proyectivo fijado por la holonomía nonádica. Sólo
después se aplican los teoremas clásicos de códigos, retículos y álgebras de
operadores de vértice a los objetos producidos.

\subsection{\texorpdfstring{Despliegue APP y cono graduado del índice \(54\)}{Despliegue APP y cono graduado del índice 54}}

Sea \(R\) la involución que intercambia las hojas APP y sean
\(P_+=(I+R)/2\), \(P_-=(I-R)/2\) sus proyectores. El bloque central
\begin{equation}
 B_c=7I+2R=9P_++5P_-
 \label{eq:exc-bloque-central-app}
\end{equation}
produce cuatro objetos de tipos distintos:
\begin{equation}
 9-5=4,\qquad (B_c)_{i\ne j}=2,\qquad
 \sum M_\Pi-\sum M_\Sigma=51-45=6,\qquad
 9\cdot6=459-405=54.
 \label{eq:exc-despliegue-42654}
\end{equation}
El \(4\) es el salto espectral entre hojas, el \(2\) el acoplamiento local,
el \(6\) la diferencia de las compresiones modulares y el \(54\) la
integración de esa diferencia sobre las nueve posiciones. Por tanto,
\(4\to2\to6\to54\) designa una jerarquía de morfismos y evaluaciones, no
una cadena de igualdades entre objetos.

Fijadas la orientación \(o_{\Sigma\Pi}\), la paridad de hoja y la
normalización triangular, la rama de Fock posee el acoplamiento
\begin{equation}
 \mathcal H_{\Sigma,\Pi}^{(m)}:
 M_{\rm bal}^{C_3,-}\longrightarrow
 \operatorname{End}\!\bigl(\mathcal F_2(\mathcal V_m)\bigr)
 \otimes o_{\Sigma\Pi},
 \label{eq:exc-acoplamiento-fock}
\end{equation}
y su traza orientada satisface
\begin{equation}
 \operatorname{Tr}_{o,\mathcal F_2}
 \bigl(\mathcal H_{\Sigma,\Pi}^{(m)}(\nu)\Gamma_2(g_m)\bigr)
 =-\frac{3m\,c(\nu)}2.
 \label{eq:exc-traza-fock}
\end{equation}
Para \(m=12\) y \(c(\delta)=-3\), vale \(54\). Independientemente, el
producto ciclotómico
\begin{equation}
 t_3(q)=q^{-1}\prod_{n\ge1}(1+q^n+q^{2n})^{-12}
 =q^{-1}-12+54q-76q^2-243q^3+\cdots
 \label{eq:exc-t3}
\end{equation}
da \([q]t_3=54\).

\begin{theorem}[Cono tipado del índice común]
\label{thm:exc-cono-54}
Una vez construidos por separado el censo APP, la rama de Fock, el producto
ciclotómico, el vecino reticular de
\cref{eq:exc-vecino-indice-tres} y el módulo lunar de
\cref{eq:exc-flm}, sus lectores verifican
\begin{equation}
 D_{\rm APP}
 =\operatorname{Tr}_{o,\mathcal F_2}
  \bigl(\mathcal H_{\Sigma,\Pi}^{(12)}(\delta)\Gamma_2(g_{12})\bigr)
 =[q]t_3=v^2=\frac{196884}{5\cdot729+1}=54.
 \label{eq:exc-cono-54}
\end{equation}
La igualdad reúne evaluaciones en un cono hacia \(\mathbb Z\); no
identifica sus dominios ni usa \(196884\) para seleccionar el corredor.
\end{theorem}

\begin{proof}
Las tres primeras igualdades son
\eqref{eq:exc-despliegue-42654}, \eqref{eq:exc-traza-fock} y
\eqref{eq:exc-t3}. Para el vector radial
\(v=4\rho_1+\rho_2+\cdots+\rho_{12}\) de
\eqref{eq:exc-vector-vecino}, la ortogonalidad de los factores \(A_2\) da
\(v^2=4^2\cdot2+11\cdot2=54\). Finalmente,
\(5\cdot729+1=3646\) y \(54\cdot3646=196884\). Cada identidad se prueba en
su propio codominio antes de compararlas.
\end{proof}

\subsection{\texorpdfstring{Forma de Paley y raíz cuadrática interna de \(-I_6\)}{Forma de Paley y raíz cuadrática interna de -I6}}

Sea
\[
 \mathcal U_9=(\mathbb Z/9\mathbb Z)^\times=\{1,2,4,8,7,5\},
\]
en el orden de las potencias de \(2\), y sea
\begin{equation}
 \theta:\mathcal U_9\xrightarrow{\sim}\mathbb P^1(\mathbb F_5),
 \qquad
 \begin{array}{c|cccccc}
 u&1&2&4&8&7&5\\ \hline
 \theta(u)&\infty&0&1&2&3&4
 \end{array}.
 \label{eq:exc-calibre-proyectivo}
\end{equation}
Esta biyección es un calibre de coordenadas del lector; no se usa como
homomorfismo. Para el carácter cuadrático
\(\chi:\mathbb F_5\to\{-1,0,1\}\), defínase
\begin{equation}
 C_P=
 \begin{pmatrix}
 0& 1& 1& 1& 1& 1\\
 1& 0& 1&-1&-1& 1\\
 1& 1& 0& 1&-1&-1\\
 1&-1& 1& 0& 1&-1\\
 1&-1&-1& 1& 0& 1\\
 1& 1&-1&-1& 1& 0
 \end{pmatrix},
 \label{eq:exc-matriz-paley}
\end{equation}
donde las entradas finitas fuera de la diagonal son \(\chi(x-y)\) y las
entradas incidentes con \(\infty\) son uno.

\begin{lemma}[Correlación cuadrática]
\label{lem:exc-correlacion-cuadratica}
Si \(a\ne b\) en \(\mathbb F_5\), entonces
\[
 \sum_{t\in\mathbb F_5}\chi(t-a)\chi(t-b)=-1.
\]
\end{lemma}

\begin{proof}
La sustitución \(u=(t-a)/(b-a)\) reduce la suma a
\(\sum_u\chi(u(u-1))\). Para \(u=0,1,2,3,4\), los sumandos son
\(0,0,1,-1,-1\).
\end{proof}

\begin{proposition}[Identidad de conferencia y reducción ternaria]
\label{prop:exc-paley-witt}
La matriz \(C_P\) es simétrica y satisface
\[
 C_PC_P^{\mathsf T}=5I_6.
\]
Su reducción \(A_W=C_P\bmod3\) verifica
\begin{equation}
 \boxed{A_W^2=-I_6\quad\text{en }M_6(\mathbb F_3).}
 \label{eq:exc-aw-cuadrado}
\end{equation}
\end{proposition}

\begin{proof}
Cada fila de \(C_P\) tiene norma cuadrada cinco. Dos filas finitas
distintas tienen producto escalar \(1-1=0\) por el
\cref{lem:exc-correlacion-cuadratica}; la suma de un carácter no trivial
prueba también la ortogonalidad con la fila de \(\infty\). Al reducir módulo
tres se obtiene
\(A_WA_W^{\mathsf T}=5I_6=2I_6=-I_6\). Como \(A_W\) es simétrica,
\eqref{eq:exc-aw-cuadrado} sigue.
\end{proof}

La identidad \eqref{eq:exc-aw-cuadrado} construye la estructura compleja
finita del corredor. En efecto,
\[
 \mathbb F_9=\mathbb F_3[\iota]/(\iota^2+1),
 \qquad
 P_{\pm\iota}=\frac12(I\mp\iota A_W)
\]
son proyectores complementarios de rango tres y dotan a
\(\mathbb F_3^6\) de estructura de \(\mathbb F_9\)-módulo de rango tres.
La acción es explícita:
\begin{equation}
 (a+b\iota)\cdot v=av+b(vA_W).
 \label{eq:exc-accion-f9}
\end{equation}
Por tanto, \(3^6=9^3\) no se usa como sustituto cardinal de una acción.

\subsection{Código ternario autodual y sistema de Witt}

Defínase
\begin{equation}
 c:\mathbb F_3^6\longrightarrow\mathbb F_3^{12},
 \qquad c(q)=(q,qA_W),
 \qquad C_W:=\operatorname{im}c.
 \label{eq:exc-codigo-grafico}
\end{equation}

\begin{theorem}[Código ternario producido por la incidencia]
\label{thm:exc-codigo-ternario}
La aplicación \(c\) es inyectiva y
\begin{equation}
 C_W=C_W^\perp,
 \qquad
 W_{C_W}(y)=1+264y^6+440y^9+24y^{12}.
 \label{eq:exc-enumerador-cw}
\end{equation}
En particular, \(C_W\) tiene parámetros \([12,6,6]_3\).
\end{theorem}

\begin{proof}
La primera proyección de \(c(q)\) es \(q\), luego \(c\) es inyectiva.
Con \(G_W=(I_6\mid A_W)\),
\[
 G_WG_W^{\mathsf T}=I_6+A_WA_W^{\mathsf T}=0
\]
en \(\mathbb F_3\). Así \(C_W\subseteq C_W^\perp\), y la igualdad de
dimensiones da la autodualidad. El enumerador es la suma finita exacta
\[
 \sum_{q\in\mathbb F_3^6}
 y^{\operatorname{wt}(q)+\operatorname{wt}(qA_W)}.
\]
Evaluar las \(3^6=729\) filas de la matriz explícita
\eqref{eq:exc-matriz-paley} da, respectivamente, los censos
\(1,264,440,24\) en pesos \(0,6,9,12\), y cero en los restantes. La suma
de los cuatro censos es \(729\), por lo que no queda ninguna palabra fuera
del certificado finito.
\end{proof}

Los \(264\) vectores de peso seis aparecen en pares \(\{x,-x\}\); sus
\(132\) soportes forman el sistema de Steiner
\begin{equation}
 W_{12}=S(5,6,12).
 \label{eq:exc-w12}
\end{equation}
Éste es el reconocimiento, después de construir \(C_W\), del código ternario
extendido de Golay y de su diseño de Witt
\citep{MacWilliamsSloane1977,ConwaySloane1999}. No se introduce el diseño
para seleccionar la matriz \(A_W\).

La prolongación binaria se mantiene separada. Si
\(\mathcal G_{24}\subset\mathbb F_2^{24}\) es el código binario extendido
de Golay, \(D\) una dodecada y
\(\ell:\{1,\ldots,12\}_{\rm HMT}\to D\) un isomorfismo de diseños, los
soportes de peso ocho forman \(W_{24}=S(5,8,24)\). Para una tétrada fija,
\[
 \lambda_4(W_{12})=4,
 \qquad
 \lambda_4(W_{24})=5,
\]
de modo que cuatro octadas elevan las cuatro hexadas incidentes y una quinta
es transversal. El dato \((\mathcal G_{24},D,\ell)\) pertenece al codominio
binario; no existe una flecha canónica \(W_{24}\to\Lambda_{24}\).

\subsection{Pegado ternario y vecino de Leech}

Sea \(Q=A_2^{12}\), con \(\det Q=3^{12}\). El discriminante de cada
factor \(A_2\) es \(\mathbb F_3\), y el pegado definido por
\eqref{eq:exc-codigo-grafico} es
\begin{equation}
 N:=\{x\in Q^*:x\bmod Q\in C_W\}.
 \label{eq:exc-niemeier-pegado}
\end{equation}

\begin{theorem}[Retículo de Niemeier producido por el código]
\label{thm:exc-niemeier}
El retículo \(N\) es positivo definido, par y unimodular de rango
veinticuatro, y su sistema de raíces es \(A_2^{12}\). Por consiguiente,
\(N\cong N(A_2^{12})\).
\end{theorem}

\begin{proof}
La autodualidad de \(C_W\) hace integral el pegado. Todos sus pesos son
múltiplos de tres, lo que hace par la extensión. Además,
\([N:Q]=|C_W|=3^6\), y por ello
\[
 \det N=\frac{\det Q}{[N:Q]^2}
 =\frac{3^{12}}{3^{12}}=1.
\]
Una clase discriminante no nula de \(A_2\) tiene norma mínima \(2/3\).
Como toda palabra no nula de \(C_W\) tiene peso al menos seis, todo vector
de pegado exterior a \(Q\) tiene norma al menos cuatro. Los vectores de
norma dos son exactamente las raíces de los doce factores \(A_2\). La
clasificación de los retículos pares unimodulares de rango veinticuatro
identifica el tipo indicado \citep{ConwaySloane1999}.
\end{proof}

Sea \(\rho_i\) la raíz máxima del factor \(A_2^{(i)}\), normalizada por
\(\rho_i^2=2\), y defínase
\begin{equation}
 v=4\rho_1+\rho_2+\cdots+\rho_{12},
 \qquad v^2=54.
 \label{eq:exc-vector-vecino}
\end{equation}
El marco ordenado del lector excepcional fija cuál de las doce fibras ocupa
la primera posición; ningún valor arquimediano interviene en esta selección.
Póngase
\begin{equation}
 \chi_v(x)=\langle x,v\rangle\bmod3,
 \quad N_0=\ker\chi_v,
 \quad N_v=N_0+\mathbb Z\frac v3.
 \label{eq:exc-vecino-indice-tres}
\end{equation}

\begin{lemma}[Estructura del vecino]
\label{lem:exc-estructura-vecino}
El carácter \(\chi_v\) es sobreyectivo, \([N:N_0]=3\),
\(v/3\in N_0^*\), y \(N_v\) es par y unimodular.
\end{lemma}

\begin{proof}
Para toda raíz de un factor \(A_2\), el producto con la raíz máxima es
\(\pm1\) o \(\pm2\); los coeficientes de \(v\) son congruentes con uno
módulo tres. Por tanto, ninguna raíz pertenece a \(N_0\) y \(\chi_v\) es
no trivial, luego sobreyectivo. La definición del núcleo da índice tres y
\(\det N_0=9\). Para \(x\in N_0\),
\(\langle x,v/3\rangle\in\mathbb Z\), mientras
\((v/3)^2=6\). Así la extensión por \(v/3\) es integral y par. Como
\([N_v:N_0]=3\),
\(\det N_v=9/3^2=1\).
\end{proof}

\begin{lemma}[Mínimo de las clases nuevas]
\label{lem:exc-minimo-clases}
Cada coordenada de una clase local que aparece en
\(N_0\pm v/3\) tiene norma mínima \(2/9\). En consecuencia,
\[
 \|x\pm v/3\|^2\ge12\cdot\frac29=\frac83
 \qquad(x\in N_0).
\]
\end{lemma}

\begin{proof}
Escribiendo una coordenada de \(A_2^*\) en la base de raíces simples, las
seis clases orientadas relevantes son los desplazamientos de
\(\pm\rho/3\) por las tres clases discriminantes. Completar cuadrados en la
forma de Gram
\(\bigl(\begin{smallmatrix}2&-1\\-1&2\end{smallmatrix}\bigr)\)
da mínimo \(2/9\) en las seis clases. La suma ortogonal de las doce
coordenadas da la desigualdad.
\end{proof}

\begin{theorem}[Construcción del retículo de Leech]
\label{thm:exc-leech}
El vecino \(N_v\) de \eqref{eq:exc-vecino-indice-tres} es isométrico al
retículo de Leech:
\begin{equation}
 \boxed{N_v\cong\Lambda_{24}.}
 \label{eq:exc-leech}
\end{equation}
\end{theorem}

\begin{proof}
Las raíces de \(N\) no pertenecen a \(N_0\), y las dos clases nuevas no
contienen vectores de norma dos por el
\cref{lem:exc-minimo-clases}. El \cref{lem:exc-estructura-vecino} prueba que
\(N_v\) es positivo definido, par, unimodular y de rango veinticuatro. La
unicidad del retículo con esas propiedades y sin raíces lo identifica con
\(\Lambda_{24}\) \citep{ConwaySloane1999}.
\end{proof}

\subsection{Isometría de orden tres, orbifold y módulo lunar}

La rotación de Coxeter de cada factor \(A_2\) induce una isometría
\(g_c\) de orden tres. La estabilidad del código \(C_W\) prueba que
\(g_cN=N\). Para el marco orientado que fija \(v\), los desplazamientos
\begin{equation}
 y_*:=\frac{g_cv-v}{3},
 \qquad
 z_*:=\frac{g_c^{-1}v-v}{3}
 \label{eq:exc-desplazamientos-vecino}
\end{equation}
pertenecen a \(N_0\): sus palabras discriminantes son un par
\(c_*,-c_*\in C_W\), y
\(\langle y_*,v\rangle=\langle z_*,v\rangle=-27\). Por tanto,
\begin{equation}
 g_cN_v=N_v.
 \label{eq:exc-preservacion-vecino}
\end{equation}
La acción no tiene puntos fijos en \(N_v\otimes\mathbb R\), pues cada
factor \(A_2\) posee los dos autovalores primitivos de orden tres.

Sea \(\widehat g_c\) la elevación estándar a la VOA reticular
\(V_{\Lambda_{24}}\). Sus dos autoespacios no triviales tienen dimensión
doce. La energía de vacío del módulo torcido es, por tanto,
\begin{equation}
 \rho_{g_c}
 =\frac1{4\cdot3^2}
 \sum_{k=1}^{2}k(3-k)\dim\mathfrak h_{(k)}
 =\frac43.
 \label{eq:exc-peso-torcido}
\end{equation}
La congruencia \(3^2\rho_{g_c}=12\equiv0\pmod3\) prueba que la elevación
es de tipo cero.

\begin{theorem}[Orbifold triádico]
\label{thm:exc-orbifold}
El orbifold cíclico de la elevación anterior satisface
\begin{equation}
 \operatorname{Orb}_{\mathbb Z/3}
 (V_{\Lambda_{24}},\widehat g_c)\cong V^\natural.
 \label{eq:exc-orbifold}
\end{equation}
El automorfismo dual pertenece a la clase \(3B\).
\end{theorem}

\begin{proof}
Se han verificado las tres hipótesis que no pueden sustituirse por una
igualdad de series: isometría fijo-libre de orden tres, preservación del
vecino y tipo cero. Los teoremas de orbifold cíclico identifican entonces
la extensión holomorfa con el módulo lunar y clasifican el automorfismo dual
\citep{ChenLamShimakura2016,AbeLamYamada2017,Carnahan2017}.
\end{proof}

La construcción de orden dos es una realización paralela. Para la elevación
\(\theta\) de la negación en \(\Lambda_{24}\),
\begin{equation}
 V^\natural=(V_{\Lambda_{24}})^+
 \oplus(V_{\Lambda_{24}}^T)^+.
 \label{eq:exc-flm}
\end{equation}
Las rutas de orden dos y tres se identifican por los teoremas de VOA; no se
serializan causalmente una dentro de la otra \citep{FLM1988}.

\begin{theorem}[Publicaciones del módulo lunar]
\label{thm:exc-moonshine}
Una vez construido \(V^\natural\), se obtienen
\begin{align}
 \operatorname{Aut}(V^\natural)&\cong\mathbb M,
 \label{eq:exc-aut-monster}\\
 \operatorname{ch}_{V^\natural}(\tau)
 &=J(\tau)=j(\tau)-744
 =q^{-1}+196884q+21493760q^2+\cdots,
 \label{eq:exc-caracter-j}\\
 T_g(\tau)&=\operatorname{Tr}_{V^\natural}
 \bigl(gq^{L_0-1}\bigr),\qquad g\in\mathbb M.
 \label{eq:exc-mckay-thompson}
\end{align}
Las series \(T_g\) son las funciones modulares de género cero del teorema
de Moonshine y satisfacen las identidades de replicabilidad.
\end{theorem}

\begin{proof}
La primera identificación y la construcción graduada proceden de la
realización FLM. La coordinación entre la acción graduada, las series de
McKay--Thompson y sus propiedades modulares es el teorema de Moonshine de
Conway--Norton y Borcherds
\citep{FLM1988,ConwayNorton1979,Borcherds1992}. En particular,
\((V^\natural)_0=\mathbb C\mathbf1\) y
\((V^\natural)_1=0\); el coeficiente \(196884=1+196883\) es la primera
descomposición no trivial en dimensiones de representaciones del Monstruo.
\end{proof}

\begin{falsifier}[Corredor excepcional]
\label{fals:exc-corredor}
La publicación queda refutada si falla cualquiera de los pasos efectivos:
\eqref{eq:exc-despliegue-42654}, la traza
\eqref{eq:exc-traza-fock}, el coeficiente \([q]t_3=54\),
\(C_PC_P^{\mathsf T}=5I_6\), \(A_W^2=-I_6\), la autodualidad o el
enumerador de \(C_W\), la paridad o unimodularidad del pegado, la ausencia
de raíces del vecino, su preservación por \(g_c\), el peso torcido
\(4/3\) o el tipo cero. También queda refutada una narración que utilice
\(W_{24}\) para construir \(\Lambda_{24}\), que confunda los dos conjuntos
de \(243\) por su cardinal, o que haga actuar al Monstruo directamente
sobre cifras y rutas. La acción de \(\mathbb M\) reside en \(V^\natural\).
\end{falsifier}
